% Generated draft: requires mathematical and source review.
\documentclass[11pt]{article}
\usepackage[margin=25mm]{geometry}
\usepackage{fontspec}
\setmainfont{Latin Modern Roman}
\usepackage{amsmath,amssymb,mathtools,unicode-math}
\setmathfont{Latin Modern Math}
\usepackage{enumitem,needspace,xurl,hyperref}
\setlength{\parindent}{0pt}
\setlength{\parskip}{5pt}
\setlength{\emergencystretch}{3em}
\newcommand{\N}{\mathbb{N}}
\newcommand{\Z}{\mathbb{Z}}
\newcommand{\Q}{\mathbb{Q}}
\newcommand{\R}{\mathbb{R}}
\newcommand{\C}{\mathbb{C}}
\newcommand{\F}{\mathbb{F}}
\newcommand{\A}{\mathbb{A}}
\newcommand{\PP}{\mathbb{P}}
\newcommand{\E}{\mathbb{E}}
\newcommand{\eps}{\varepsilon}
\newcommand{\dd}{\,\mathrm{d}}
\begin{document}
\section*{UnsolvedMath ID: 801079}
\textbf{Source problem code:} catalogue entry
\subsection{Definitions and mathematical statement}
The source asks about two generalized H\'enon maps \(f_1,f_2\). In the standard complex setting, an elementary generalized H\'enon map has the form
\[
(x,y)\longmapsto (y,p(y)-ax),\qquad p\in\C[y],\quad \deg p\geq 2,\quad a\in\C\setminus\{0\}.
\]
The supplied record does not specify whether its class of generalized H\'enon maps includes finite compositions or other coordinate conventions; this example does not impose an additional restriction on the source question.

For a map \(f_i\colon\C^2\to\C^2\) in the intended class, define its forward filled set by
\[
K_i^+=\left\{z\in\C^2:\sup_{n\geq 0}\|f_i^n(z)\|<\infty\right\}.
\]
Let \(W_i\) be an unstable manifold for \(f_i\), and let \(U_i\subseteq W_i\) be relatively open. The associated local unstable slice is
\[
S_i=K_i^+\cap U_i.
\]
The source does not specify the invariant points or other invariant objects whose unstable manifolds are allowed.

\textbf{Question.} Characterize the circumstances under which there exist choices of unstable manifolds \(W_1,W_2\) and relatively open subsets \(U_1\subseteq W_1\), \(U_2\subseteq W_2\) such that the slices \(S_1\) and \(S_2\) are biholomorphically equivalent.

The meaning of this equivalence requires clarification because the slices need not themselves be open complex manifolds. If the intended meaning is equivalence induced by a biholomorphism of the chosen ambient open subsets, the condition is
\[
\exists\,h\colon U_1\longrightarrow U_2
\quad\text{biholomorphic, with}\quad
h(K_1^+\cap U_1)=K_2^+\cap U_2.
\]
This is a conditional interpretation, not a qualification explicitly supplied by the source. Equivalence on smaller neighborhoods or at specified germs could give a different formulation. The record also does not require the slices to be nonempty or specify a distinguished point that they must contain. These unresolved choices must be settled before a unique precise formulation can be asserted.
\subsection{Short English statement}
Berger and Bedford ask when a local unstable slice of one generalized H\'enon map can be biholomorphically equivalent to a local unstable slice of another. A slice is the intersection of the forward bounded-orbit set with a relatively open part of an unstable manifold. The allowed unstable manifolds, nonemptiness requirements, and precise meaning of biholomorphic equivalence remain unspecified in the supplied record.
\subsection{Sources}
\begin{itemize}
\item Local UnsolvedMath registry: permanent ID \(801079\), exact source problem code \texttt{catalogue entry}, title ``Henon Maps: Open Problems --- Question 31.'' This draft uses the supplied complete source record.
\item Berger, Bedford, \emph{Henon Maps: A List of Open Problems} (2023), Question 31. Supplied source URL: \url{https://arxiv.org/abs/2312.03907}. The registry attributes the question to Berger and Bedford, with proposed year 2023.
\item Supplied provenance: extraction method \texttt{arxiv\_latex}; recorded access date 2026-10-16. This access information is supplied provenance and does not indicate independent source consultation for this draft.
\item The registry labels the problem \texttt{open}; that label is supplied curation, not an independently verified mathematical or literature assessment.
\end{itemize}
\end{document}
